% =============================================================================
% ENGLISH VERSION – Two-column (revtex4-2, APS/PRE style)
% Manual translation of borrador_rebotes_gravedad.tex (2026-07-02).
% If the Spanish master changes, update this file accordingly.
% Compile: pdflatex borrador_rbf_en.tex  (two passes)
% =============================================================================

\documentclass[aps,pre,twocolumn,reprint,superscriptaddress]{revtex4-2}

\usepackage[utf8]{inputenc}
\usepackage[T1]{fontenc}
\usepackage{amsmath,amssymb}
\usepackage{graphicx}
\usepackage{booktabs}
\usepackage{hyperref}

% --- Physics macros ----------------------------------------------------------
\newcommand{\cor}{e}
\newcommand{\coR}[1]{e_{#1}}
\newcommand{\Tstar}{T^*}
\newcommand{\eps}{\varepsilon}
\newcommand{\epsps}{\varepsilon_{ps}}
\newcommand{\gstar}{g^*}
\newcommand{\lammps}{\textsc{Lammps}}
\newcommand{\TODO}[1]{\textbf{[TODO: #1]}}

\begin{document}

\title{Energy partition in the impact of Lennard-Jones nanoclusters on FCC
surfaces: four channels, projectile-size scaling, and adhesion}

% email alternative: felipe.wasaff@uchile.cl
\author{Felipe González-Wasaff}
\email{felipe.wasaff.g@mail.pucv.cl}
\affiliation{Instituto de Matemáticas, Pontificia Universidad Católica de
Valparaíso, Valparaíso, Chile}

\author{Gonzalo Gutiérrez}
\affiliation{Departamento de Física, Facultad de Ciencias,
Universidad de Chile, Santiago, Chile}

\date{\today}

\begin{abstract}
We study by molecular dynamics (\lammps, LJ units) the energy partition in
the first bounce of a Lennard-Jones nanocluster on a crystalline FCC surface
under a constant gravitational field, reproducing and extending the work of
Gutiérrez et al.\ (2012). The canonical system is a pre-minimized FCC
cluster of $N=79$ atoms (5th complete shell, $r < 2.45\,\sigma$) on a
surface of $8\times8\times6$ unit cells ($a_0=1.55\,\sigma$) with a purely
repulsive ball--surface interaction ($c_{ps}=0$, $T=0$). We identify four
energy channels whose sum is exactly equal to the impact energy $E_0$:
\[
    \alpha_{\mathrm{cm}} + \alpha_{T_B} + \alpha_{\Delta V_B}
    + \alpha_{\mathrm{surf}} = 1,
\]
where $\alpha_{T_B}\approx0.22$ is the internal kinetic fraction of the ball
(constant for $v>0.9\,\sigma/\tau$), $\alpha_{\Delta V_B}\approx0.21$--$0.33$
is the net change of the intra-ball PE (post-impact anharmonic expansion),
and $\alpha_{\mathrm{surf}}\approx0.07$--$0.16$ is the direct excitation of
surface phonons. This breakdown complements the original analysis, which
characterized the collisional energy $U_c$ \emph{during} contact: the
partition at the first apex separates for the first time the anharmonic
expansion $\alpha_{\Delta V_B}$ from the surface excitation, and shows that
most of the non-recovered energy remains in the ball, not in the surface.

We extend the study to $N\in\{1,2,7,13,55,79,87\}$ with repulsive
interaction. The total internal fraction
$\alpha_{\mathrm{ball}}=\alpha_{T_B}+\alpha_{\Delta V_B}$ follows the
equipartition law $\alpha_{\mathrm{ball}}=(3N-3)/(3N-3+3N_c)$ with
$N_c\approx N$ for $N\geq55$, predicting the asymptotic limit
$\alpha_{\mathrm{ball}}\to\tfrac{1}{2}$ as $N\to\infty$.
For $N=13$ (smallest complete FCC shell), the ball is too stiff to reach
the equipartition plateau within the studied range $v\leq2.0\,\sigma/\tau$.

In long simulations ($5000\,\tau$, $N=79$, $h_0=28\,\sigma$), the first
bounce height is well approximated by the scale
$h_1\approx T_B/(Mg)=8.3\,\sigma$ proposed here (ratio $=0.96$), but the
system does not converge to that attractor: the CM height decays faster
than $T_B$, and after 23 bounces ($t\approx660\,\tau$) the ball remains
trapped over the surface with $h_n\lesssim1\,\sigma$.

With the full LJ interaction ($c_{ps}=1$, adhesion), the ball sticks in all
$4N\times6h_0=24$ studied cases ($v$ up to $2.0\,\sigma/\tau$). The adhesion
energy $E_{\rm adh}\propto N^{0.73}$ (contact area) exceeds in every case
the CM energy available after the bounce; an order-of-magnitude
extrapolation places the threshold at $h_0^*(N)\sim260$--$440\,\sigma$ with
$h_0^*\propto N^{-0.27}$.
\end{abstract}

\keywords{molecular dynamics, LJ nanoclusters, coefficient of restitution,
energy partition, anharmonic expansion, surface phonons, adhesion}

\maketitle

% =============================================================================
\section{Introduction}
\label{sec:intro}
% =============================================================================

The bouncing ball is one of the most studied archetypes of dissipative
physics. At the macroscopic scale, the coefficient of restitution
$\cor = |v_f|/|v_i|$ is a material property and the dissipated energy goes
entirely to the surface (or to the medium). At the nanometric scale this
assumption fails: a cluster of $N$ atoms has $3N-3$ internal degrees of
freedom ---$3N-6$ vibrational plus 3 rotational--- that can absorb energy
during each impact.

The work of Gutiérrez et al.~\cite{Gutierrez2012} was the first to study
this \emph{dissipation cascade} by molecular dynamics: an FCC argon cluster
of $N\approx79$ atoms under reduced gravity
$\gstar=0.026\,\sigma/\tau^2$, with a purely repulsive ball--surface
interaction. Their main results were: (i)~the sequence of heights $h_n$
decreases faster than the geometric series; (ii)~during the collision, the
ball--surface collisional energy $U_c$ transiently stores most of the
``dissipated'' energy, more than the vibrational energy transferred to the
surface; (iii)~for large drop heights the ball melts during the impact
(a single peak in $g(r)$ for $h_0>41.5\,\sigma$).
However, their energy decomposition was evaluated \emph{during} contact
(where $U_c$ is a transient that vanishes once the ball separates): it did
not quantify the partition at the first apex, did not separate the change
of the intra-ball PE from the direct excitation of the surface, and did not
study the dependence on $N$.

In this work we reproduce that experiment with \lammps{} using a cluster of
exactly $N=79$ atoms (complete pre-minimized FCC shell, not created with
\texttt{create\_atoms}) and we extend the analysis in four directions:
\textbf{(i)}~decomposition into \emph{four} independently measured channels;
\textbf{(ii)}~sweep over $N\in\{1,2,7,13,55,79,87\}$ to calibrate the
equipartition formula and determine its scaling;
\textbf{(iii)}~long simulation ($5000\,\tau$) to test the prediction
$h_1\approx T_B/(Mg)$ and the multi-bounce behavior;
\textbf{(iv)}~sweep with adhesion ($c_{ps}=1$) to determine the threshold
$h_0^*(N)$.

% =============================================================================
\section{Model and definitions}
\label{sec:modelo}
% =============================================================================

\subsection{Hamiltonian and energy decomposition}
\label{sec:hamiltoniano}

The system consists of a cluster of $N$ atoms (``projectile'') and an FCC
surface of $N_s$ atoms. The LJ potential between any pair at distance $r$ is
\begin{equation}
    U_{\mathrm{LJ}}(r) = 4\eps \left[ \left(\frac{\sigma}{r}\right)^{12}
    - c_{ij}\left(\frac{\sigma}{r}\right)^{6} \right],
    \quad r < r_c = 2.5\,\sigma,
    \label{eq:lj}
\end{equation}
with $c_{ij}=1$ for ball--ball and surface--surface pairs, and $c_{ij}=0$
(purely repulsive, $c_{ps}=0$) or $c_{ij}=1$ (adhesion, $c_{ps}=1$) for
ball--surface pairs.

The internal energy of the ball decomposes as
\begin{equation}
    E_{\mathrm{int}} =
    \underbrace{\sum_{i=1}^{N}\tfrac{1}{2}m|\mathbf{v}_i-\mathbf{V}_{\rm cm}|^2}_{T_B}
    +
    \underbrace{\sum_{i<j}U_{\mathrm{LJ}}(r_{ij})}_{V_B}.
    \label{eq:TB_VB}
\end{equation}
During free flight $T_B$ and $V_B$ exchange energy (internal oscillations)
while $E_{\rm int}=T_B+V_B$ remains constant.

\subsection{Energy balance of the first bounce}
\label{sec:balance}

We denote by $E_0 = Mg h_0^{\rm ef}$ the initial mechanical energy, where
$h_0^{\rm ef}=Z_{\rm cm}(0)-z_0$ is the effective height and $z_0$ is the
mean position of the minima of $Z_{\rm cm}$ (contact level). Measuring each
quantity at the first apex (superscript $(1)$) and immediately before the
impact (reference, superscript $(0)$), we define the four channels:
\begin{align}
    \alpha_{\rm cm}         &= h_1/h_0^{\rm ef} = e_1^2, \label{eq:acm}\\
    \alpha_{T_B}            &= T_B^{(1)}/E_0,   \label{eq:atb}\\
    \alpha_{\rm surf}       &= \Delta(K_S + V_{\rm surf})/E_0, \label{eq:asurf}\\
    \alpha_{\Delta V_B}     &= 1-\alpha_{\rm cm}-\alpha_{T_B}-\alpha_{\rm surf}.
    \label{eq:avb}
\end{align}
The fourth fraction $\alpha_{\Delta V_B}=(V_B^{(1)}-V_B^{(0)})/E_0$ can be
obtained in two ways: by subtraction (Eq.~\eqref{eq:avb}, which guarantees
$\sum\alpha_i=1$ exactly) or by measuring $V_B$ directly in \lammps{}.
The difference between the two defines the \emph{closure error}
$\varepsilon_{\rm closure}=|\sum\alpha_i^{\rm measured}-1|$, which we report
in Section~\ref{sec:lammps} as a quality check of the partition.

From the observed plateau $\alpha_{T_B}\approx0.22$
(Section~\ref{sec:4canales}), we propose the height scale for the first
bounce
\begin{equation}
    h^* = T_B^{(1)}/(Mg),
    \label{eq:hstar}
\end{equation}
which predicts $h_1\approx\alpha_{T_B}\,h_0^{\rm ef}$ when
$\alpha_{\rm cm}\approx\alpha_{T_B}$.

\subsection{Equipartition formula}
\label{sec:equiparticion}

We propose that the internal fraction of the first bounce follows an
equipartition law among the degrees of freedom active during contact:
\begin{equation}
    \alpha_{\rm int}(N) = \frac{3N-3}{3N-3+3N_c},
    \label{eq:equiparticion}
\end{equation}
where $N_c$ is the effective number of oscillators of the contact
footprint. The numerator counts the $3N-3$ internal degrees of freedom of
the ball (vibrational plus rotational): a not perfectly central impact also
excites rotation, so we do not subtract the 3 rotational modes. For
$N\gtrsim13$ the difference between $3N-3$ and $3N-6$ is smaller than the
remaining uncertainties; for $N=2$ and $N=7$ it is not, and those sizes are
treated only as qualitative controls.
If $\alpha_{\rm int}$ is identified with $\alpha_{T_B}$ alone, one obtains
$N_c\approx290$ (with no obvious geometric interpretation). If it is
identified with the total internal fraction
$\alpha_{\rm ball}=\alpha_{T_B}+\alpha_{\Delta V_B}$ and calibrated with the
data for $N\geq55$, one obtains $N_c\approx N$, which predicts the
asymptotic limit
\begin{equation}
    \alpha_{\rm ball}(N)\xrightarrow{N\to\infty}\frac{1}{2}.
    \label{eq:limite}
\end{equation}

% =============================================================================
\section{Computational details}
\label{sec:lammps}
% =============================================================================

The simulations are performed with \lammps~7~Feb~2024~\cite{Thompson2022}.

\paragraph{Surface.}
FCC lattice of $8\times8\times6$ unit cells ($a_0=1.5496\,\sigma$),
periodic boundary conditions in $xy$, non-periodic in $z$.
The 2 bottom layers are frozen (\texttt{setforce 0 0 0}).
No thermostat (global NVE, initial $T=0$).

\paragraph{Cluster.}
The $N$-atom cluster is built by selecting the $N$ positions closest to the
origin of an FCC lattice with the same lattice parameter. Complete
spherical shells correspond to $N\in\{1,13,19,43,55,79,87,\ldots\}$;
of these we simulate $N\in\{1,13,55,79,87\}$, and we add $N=2$ and $N=7$
(incomplete shells) as controls of the molecular regime.
Each cluster is minimized in vacuum (CG + FIRE algorithms) down to a
maximum force $<10^{-6}\,\eps/\sigma$, yielding cohesion energies per atom
of $-5.24\,\eps$ ($N=79$) and $-5.19\,\eps$ ($N=87$). This protocol
guarantees exactly $N$ atoms for all initial heights, unlike the
\texttt{create\_atoms region sphere} method of Gutiérrez et al.\ (2012),
which produces a variable $N$ (79--84) depending on the position of the
center. The ball is placed with zero CM velocity at $h_0$ above the box
level.

\paragraph{Gravity and bounce detector.}
Gravity $\gstar=0.026\,\sigma/\tau^2$ is applied only to the ball atoms
(\texttt{fix gravity}). $Z_{\rm cm}(t)$ is recorded every 20 steps
($\Delta t_{\rm print}=0.1\,\tau$). Minima and maxima of $Z_{\rm cm}$ are
identified in post-analysis with \texttt{scipy.signal.find\_peaks}.

\paragraph{Ball--surface interaction.}
For $c_{ps}=0$: pre-generated repulsive LJ table (pure $4/r^{12}$, with
analytic force $F=-dV/dr$), tabulated over $r\in[0.75,2.5]\,\sigma$ with
5000 points and linear interpolation (\texttt{pair\_style table}).
For $c_{ps}=1$: full LJ with \texttt{pair\_coeff 1 2 lj/cut 1.0 1.0 2.5}.
The switch is controlled by the \lammps{} variable \texttt{c\_ps} in
\texttt{in.bounce\_cluster}.

\paragraph{Integration and energy conservation.}
The equations of motion are integrated with velocity-Verlet
(\texttt{fix nve}) and time step $\Delta t = 5\times10^{-3}\,\tau$
(the original study~\cite{Gutierrez2012} used Beeman with
$\Delta t = 4.65\times10^{-4}\,\tau$). All potentials are truncated
\emph{and shifted} at $r_c=2.5\,\sigma$ (\texttt{pair\_modify shift yes});
the shift does not alter forces or trajectories, but it is essential for
the energy bookkeeping: without it, every pair crossing $r_c$ introduces a
jump of $\pm|U(r_c)|=0.016\,\eps$ in the recorded potential energy, and the
accumulated error reaches 2--6\% of $E_0$ per impact. With shifted
potentials, the total energy (including the gravitational PE) is conserved
to $|\Delta E|<0.03\,\eps$ ($\sim4\times10^{-4}\,E_0$) over $120\,\tau$
including a full impact, which validates $\Delta t=5\times10^{-3}\,\tau$.
As a check of the partition, we measure $\alpha_{\Delta V_B}$ independently
(directly from $V_B$) in addition to the subtraction of
Eq.~\eqref{eq:avb}: the resulting closure error is
$|\sum\alpha_i-1|<10^{-4}$ for the first bounce of $N=79$,
$h_0=28.15\,\sigma$.

% =============================================================================
\section{Results}
\label{sec:resultados}
% =============================================================================

% -----------------------------------------------------------------------------
\subsection{Cross validation: reproduction of Gutiérrez et al.\ (2012)}
\label{sec:reproduccion}
% -----------------------------------------------------------------------------

We reproduce the original parameters ($c_{ps}=0$, $T=0$, $\gstar=0.026$,
6 heights from $h_0=11.03$ to $79.53\,\sigma$) with our pre-minimized
cluster of $N=79$ atoms. The effective contact level is
$z_0\approx11.5\,\sigma$ (CM height at the minima).

The COR of the first bounce ranges from $e_1=0.92$
($h_0^{\rm ef}=8.8\,\sigma$) to $e_1=0.59$ ($h_0^{\rm ef}=79.0\,\sigma$).
Reference~\cite{Gutierrez2012} does not report numerical values of $e_1$,
but it does report the height sequences $h_n$ (their Fig.~7), which we
digitized for the quantitative comparison of
Table~\ref{tab:comparacion2012}. The digitization also reveals a difference
of convention: in~\cite{Gutierrez2012} the parameter $h_0$ is the
\emph{absolute} CM height (the curves of their Fig.~7 start exactly at
$h_0$ and converge to the rest level $z_{\rm rest}\approx10.3\,\sigma$),
whereas here $h_0$ is measured from the top of the surface box. The
comparison must therefore be made at comparable effective drop height
$h_0^{\rm ef}$, not at equal nominal $h_0$.

\begin{table*}[tp]
\centering
\small
\caption{Comparison with Fig.~7 of~\cite{Gutierrez2012} (digitized values;
uncertainty $\approx\pm0.7\,\sigma$ in the heights, which dominates the
error of $e_1$). The 2012 effective heights are
$h_0^{\rm ef}=h_0-z_{\rm rest}$ with $z_{\rm rest}\approx10.3\,\sigma$.
The 2012 row $h_0=11.03$ is a drop of $<1\,\sigma$ (trivially
quasi-elastic) and does not admit a reliable $e_1$.}
\label{tab:comparacion2012}
\begin{tabular}{c|ccc|cc}
\toprule
 & \multicolumn{3}{c|}{Gutiérrez et al.\ 2012 (Fig.~7)} &
 \multicolumn{2}{c}{This work ($N=79$)} \\
$h_0$ ($\sigma$) & $h_0^{\rm ef}$ & $h_1^{\rm ef}$ & $e_1$ &
$h_0^{\rm ef}$ & $e_1$ \\
\midrule
11.03 &  0.6 &  0.4 & ---  &  8.81 & 0.924 \\
18.64 &  8.4 &  6.8 & 0.90 & 17.21 & 0.708 \\
28.15 & 17.9 &  3.3 & 0.43 & 25.91 & 0.552 \\
41.47 & 31.3 &  8.2 & 0.51 & 41.80 & 0.633 \\
58.60 & 48.4 & 14.3 & 0.54 & 56.93 & 0.614 \\
79.53 & 69.2 & 18.2 & 0.51 & 79.03 & 0.587 \\
\bottomrule
\end{tabular}
\end{table*}

Comparing at similar $h_0^{\rm ef}$, the agreement is excellent in the
quasi-elastic regime ($e_1=0.90$ in 2012 at $h_0^{\rm ef}=8.4\,\sigma$
versus $0.92$ here at $8.8\,\sigma$) and good at high energy, where the
2012 values lie systematically $\sim0.06$--$0.09$ below (e.g.\ $0.51$ at
$h_0^{\rm ef}=69\,\sigma$ versus $0.59$ at $79\,\sigma$). That difference
is attributable to the 2012 cluster being generated with
\texttt{create\_atoms} without pre-minimization (variable $N$ of 79--84 and
residual internal energy, which opens additional absorption channels) and
to the different integrator (Beeman in 2012, velocity-Verlet here), so only
statistical agreement can be expected. The 2012 case
$h_0^{\rm ef}=17.9\,\sigma$ ($e_1=0.43$) is an outlier with respect to the
trend of both data sets; note that our series also shows non-monotonicity
in that region (minimum $e_1=0.55$ at $h_0^{\rm ef}=25.9\,\sigma$,
Table~\ref{tab:particion4}), which suggests genuine sensitivity to the
internal phase of the ball at impact. The faster-than-geometric decay of
the heights ($e_1\neq\mathrm{const}$) confirms the dissipation cascade in
both data sets.

% -----------------------------------------------------------------------------
\subsection{Four-channel partition for \texorpdfstring{$N=79$}{N=79}}
\label{sec:4canales}
% -----------------------------------------------------------------------------

Table~\ref{tab:particion4} shows the decomposition of the first bounce for
$N=79$, $c_{ps}=0$, $T=0$. The central result is that for
$v_{\rm imp}\geq0.9\,\sigma/\tau$, the anharmonic expansion greatly exceeds
the direct surface excitation: $\alpha_{\Delta V_B}\approx21$--$33\%$
versus $\alpha_{\rm surf}\approx7$--$16\%$.

\begin{table*}[tp]
\centering
\small
\caption{Energy partition of the first bounce into four channels for
$N=79$ (pre-minimized FCC cluster), $c_{ps}=0$, $T=0$.
$\sum\alpha_i=1$ by construction.}
\label{tab:particion4}
\begin{tabular}{ccccccc}
\toprule
$h_0^{\rm ef}$ & $v_{\rm imp}$ & $e_1$ &
$\alpha_{\rm cm}$ & $\alpha_{T_B}$ & $\alpha_{\Delta V_B}$ &
$\alpha_{\rm surf}$ \\
($\sigma$) & ($\sigma/\tau$) & & & & & \\
\midrule
 8.81 & 0.677 & 0.924 & 0.853 & 0.021 & 0.056 & 0.070 \\
17.21 & 0.946 & 0.708 & 0.502 & 0.213 & 0.208 & 0.077 \\
25.91 & 1.161 & 0.552 & 0.304 & 0.319 & 0.310 & 0.066 \\
41.80 & 1.474 & 0.633 & 0.401 & 0.257 & 0.264 & 0.078 \\
56.93 & 1.721 & 0.614 & 0.378 & 0.197 & 0.328 & 0.098 \\
79.03 & 2.027 & 0.587 & 0.345 & 0.204 & 0.294 & 0.156 \\
\bottomrule
\end{tabular}
\end{table*}

The internal kinetic fraction $\alpha_{T_B}$ is practically constant for
$v_{\rm imp}\geq0.9\,\sigma/\tau$ ($\alpha_{T_B}\approx0.21$--$0.24$).
This plateau quantifies, at the first apex, the internal heating of the
ball that~\cite{Gutierrez2012} described qualitatively during contact.
The total internal fraction
$\alpha_{\rm ball}=\alpha_{T_B}+\alpha_{\Delta V_B}$ varies between 0.42
and 0.63, exceeding $\alpha_{\rm cm}$ in all cases for
$v>1.0\,\sigma/\tau$.

The direct surface excitation $\alpha_{\rm surf}$ is the smallest
contribution in the system, between 7 and 16\% of $E_0$.
This result refines the picture suggested by the during-contact analysis
of~\cite{Gutierrez2012}: there, the collisional energy $U_c$ dominates the
\emph{transient} balance, but $U_c$ vanishes once the ball separates. The
partition at the apex shows that most of that energy returns to the ball as
anharmonic expansion ($\alpha_{\Delta V_B}$) and does not remain in the
surface. Both descriptions are compatible: they measure the distribution at
different instants of the bounce.

Figure~\ref{fig:secuencia} illustrates this sequence with simulation
snapshots ($h_0=28.15\,\sigma$): at maximum compression the energy is
mostly stored as intra-ball and contact PE (transient), while at the first
apex the ball retains $\sim$62\% of $E_0$ as internal energy
($T_B+\Delta V_B$) ---visible as the hot ``speckling'' of its atoms--- and
the surface only $\sim$6\%.

\begin{figure*}[tp]
    \centering
    \includegraphics[width=\textwidth]{../data/outputs/fig_sequence_en.png}
    \caption{Four instants of the first bounce ($N=79$, $h_0=28.15\,\sigma$,
    $c_{ps}=0$): \textbf{(a)}~free fall, \textbf{(b)}~maximum compression,
    \textbf{(c)}~post-impact expansion and \textbf{(d)}~first apex.
    Ball atoms are colored by their internal KE (CM motion removed, scale
    $0$--$0.30\,\eps$) and surface atoms by their KE ($0$--$0.01\,\eps$,
    a $\times30$ more sensitive scale to visualize the phonons). The bar
    below each panel shows the instantaneous partition of $E_0$; $U_c$
    (gray) is appreciable only during contact, and at the apex the bar
    reproduces the four channels of Table~\ref{tab:particion4} (with
    $U_F=\alpha_{\rm cm}E_0$).}
    \label{fig:secuencia}
\end{figure*}

% -----------------------------------------------------------------------------
\subsection{Scaling with \texorpdfstring{$N$}{N}: calibration of \texorpdfstring{$N_c$}{Nc}}
\label{sec:sweep_N}
% -----------------------------------------------------------------------------

We run the same six heights with $N\in\{1,2,7,13,55,79,87\}$, using
pre-minimized clusters in complete FCC shells and averaging over the points
with $v_{\rm imp}>1.1\,\sigma/\tau$. The results are shown in
Table~\ref{tab:sweep_N} and Figure~\ref{fig:sweep_N}.

\begin{table*}[tp]
\centering
\caption{First-bounce fractions averaged over $v_{\rm imp}>1.1\,\sigma/\tau$,
as a function of $N$. The last two columns show the prediction of
Eq.~\eqref{eq:equiparticion} calibrated with $\alpha_{T_B}$ ($N_c=290$) and
with $\alpha_{\rm ball}$ ($N_c=N$).}
\label{tab:sweep_N}
\begin{tabular}{cccccccc}
\toprule
$N$ & shells & $\langle e_1\rangle$ & $\alpha_{T_B}$ & $\alpha_{\rm ball}$ &
$\alpha_{\rm surf}$ &
$\hat\alpha_{T_B}(N_c{=}290)$ & $\hat\alpha_{\rm ball}(N_c{=}N)$ \\
\midrule
 1 & 1 (complete) & — & 0.000 & 0.000 & —     & 0.000 & 0.000 \\
 2 & 1+1          & 0.902 & 0.002 & 0.003 & 0.183 & 0.003 & 0.333 \\
 7 & 1+6/12$^a$   & 0.840 & 0.034 & 0.093 & 0.198 & 0.020 & 0.462 \\
13$^b$ & 2 (complete) & 0.735 & 0.066 & 0.107 & 0.339 & 0.040 & 0.480 \\
55 & 4 (complete) & 0.521 & 0.217 & 0.474 & 0.254 & 0.157 & 0.495 \\
79 & 5 (complete) & 0.597 & 0.244 & 0.543 & 0.100 & 0.212 & 0.497 \\
87 & 6 (complete) & 0.594 & 0.250 & 0.543 & 0.103 & 0.229 & 0.497 \\
\bottomrule
\end{tabular}\\
\footnotesize $^a$ $N=7$ is an incomplete shell (7 of 13 possible in the
2nd FCC shell); its values show larger scatter.\\
\footnotesize $^b$ $N=13$ does not reach the plateau within
$v\leq2.0\,\sigma/\tau$ (see text); its average is not representative and
is excluded from the calibration of $N_c$.\\
\footnotesize For $N=1$ there are no internal modes
($\alpha_{T_B}\equiv0$); that row is a control. For $N=2$, almost all the
non-recovered energy goes to the surface ($\alpha_{\rm surf}=0.183$, with
$\alpha_{\rm cm}=\langle e_1\rangle^2\approx0.81$): the dimer barely
couples to its 3 internal modes in vertical impacts.
\end{table*}

\paragraph{Calibration of $N_c$.}
For $\alpha_{T_B}$ alone (Eq.~\eqref{eq:equiparticion}), the fit gives
$N_c^{\rm eff}\in\{169,195,241,258\}$ for $N\in\{13,55,79,87\}$, growing
weakly with $N$ ($N_c\propto N^{0.20}$) with no direct geometric
interpretation. In contrast, using $\alpha_{\rm ball}$ in
Eq.~\eqref{eq:equiparticion} one obtains $N_c\approx N$ for $N\geq55$
(column $\hat\alpha_{\rm ball}$). The agreement is not perfect: the data
($\alpha_{\rm ball}=0.543$ for $N=79$ and $87$) systematically exceed the
$N_c=N$ prediction (0.497) and even the asymptotic bound $1/2$, which
strictly implies $N_c<N$: inverting Eq.~\eqref{eq:equiparticion} gives
$N_c/N=1.09$, $0.83$ and $0.83$ for $N=55$, $79$ and $87$. The systematic
excess may originate in the rotational energy of the ball (included in
$\alpha_{\rm ball}$ but not in the harmonic-oscillator counting) or in an
underestimation of $\alpha_{\rm surf}$; discriminating between the two
requires measuring the post-bounce angular momentum
(\texttt{compute angmom}), left for a revision. With that caveat, the
identification $N_c\sim N$ ---each ball atom effectively coupled to one
surface oscillator--- captures the trend and predicts the limit
$\alpha_{\rm ball}\to1/2$ as $N\to\infty$.

\paragraph{Small-ball regime ($N=13$).}
The $N=13$ cluster (the smallest complete spherical FCC shell, radius
$1.10\,\sigma$) does not reach the equipartition plateau within
$v\leq2.0\,\sigma/\tau$: $\alpha_{T_B}$ keeps growing from $0.003$ (low
energy) to $0.152$ (highest studied velocity), while the equipartition
prediction would give $\alpha_{T_B}\approx0.04$ ($N_c=290$) or
$\alpha_{\rm ball}\approx0.48$ ($N_c=N$). The stiffness of the single-shell
cluster prevents the fast thermalization that enables the plateau in larger
clusters.

\begin{figure*}[tp]
    \centering
    \includegraphics[width=\textwidth]{../data/outputs/fig_alpha_TB_full_sweep_en.pdf}
    \caption{Sweep over $N\in\{1,2,7,13,55,79,87\}$, $c_{ps}=0$, $T=0$.
    \textbf{(a)}~$\alpha_{T_B}$ vs $v_{\rm imp}$ for each $N$; the vertical
    dotted line marks $v^*=1.1\,\sigma/\tau$ (onset of the plateau).
    \textbf{(b)}~$\langle\alpha_{T_B}\rangle_{\rm plateau}$ vs $N$ (error
    bars are the standard deviation over the points with $v>v^*$), compared
    with Eq.~\eqref{eq:equiparticion} for $N_c=290$ (solid blue curve) and
    with the fit $N_c(N)=99\,N^{0.20}$ (dashed black curve). The open
    symbol of $N=13$ indicates that this size does not reach the plateau.
    The inset shows $N_c^{\rm eff}(N)$ and its power-law fit.}
    \label{fig:sweep_N}
\end{figure*}

% -----------------------------------------------------------------------------
\subsection{Long simulation: testing the prediction \texorpdfstring{$h^*=T_B/(Mg)$}{h*}}
\label{sec:largo}
% -----------------------------------------------------------------------------

To test the scale $h^*=T_B/(Mg)$ proposed in Section~\ref{sec:balance}
---and in particular whether the bounce height \emph{converges} to it--- we
run a $5000\,\tau$ simulation ($N=79$, $h_0=28.15\,\sigma$, $c_{ps}=0$,
same parameters).

\paragraph{First bounce.}
The measured height $h_1=7.91\,\sigma$ is in excellent agreement with the
prediction $h^*(T_B^{(1)})=T_B^{(1)}/(Mg)=8.27\,\sigma$ (ratio $=0.96$).
This result confirms that the relation $h_1\approx T_B/(Mg)$ is valid for
the first impact.

\paragraph{Successive bounces and collapse.}
Contrary to what the name ``$h^*$'' suggests as a stationary attractor, the
ratio $h_n/h^*(T_B^{(n)})$ decreases monotonically bounce after bounce
(from 0.96 at $n=1$ to $\sim0.18$ at $n=20$--$23$). The CM height $h_n$
decays geometrically ($e_1^{\rm eff}\approx0.97$ per bounce), while $T_B$
decreases more slowly (from $17.0\,\eps$ at $n=1$ to $11.8\,\eps$ at
$n=23$, a factor $\times0.7$).

After 23 bounces ($t\approx660\,\tau$), the apex height has dropped to
$h_n\approx1.0\,\sigma$: the ball no longer effectively separates from the
surface. At $t\approx713\,\tau$ the CM falls below $z_0$ (collapse): ball
atoms leak through the gaps of the surface lattice since the purely
repulsive wall has no upper barrier in $+z$.

\paragraph{Interpretation.}
The prediction $h_1\approx h^*=T_B/(Mg)$ is a consequence of the first
impact ($\alpha_{\rm cm}\approx\alpha_{T_B}$ in the first bounce, see
Table~\ref{tab:particion4}), not a dynamical attractor. In successive
bounces, the additional available CM energy (falling from $h_n\ll h_0$) is
small, so $T_B$ grows little while $h_n$ keeps decreasing through
inelasticity.

\begin{figure*}[tp]
    \centering
    \includegraphics[width=\textwidth]{../data/outputs/fig_paso_B_convergencia_en.pdf}
    \caption{Long simulation ($5000\,\tau$, $N=79$, $h_0=28.15\,\sigma$,
    $c_{ps}=0$).
    \textbf{(a)}~Trajectory $Z_{\rm cm}(t)-z_0$ during the first
    $1000\,\tau$; the dashed red line marks the collapse through the
    lattice at $t\approx713\,\tau$ (an artifact of the purely repulsive
    potential). The inset shows the full $5000\,\tau$.
    \textbf{(b)}~Apex heights $h_n$ (blue circles) and prediction
    $h^*_n=T_B^{(n)}/(Mg)$ (red squares) per bounce number.
    \textbf{(c)}~Ratio $h_n/h^*_n$: the first-bounce value (0.96) is the
    only one that approaches unity (dashed line).}
    \label{fig:largo}
\end{figure*}

% -----------------------------------------------------------------------------
\subsection{Adhesion: sweep with \texorpdfstring{$c_{ps}=1$}{c\_ps=1}}
\label{sec:adhesion}
% -----------------------------------------------------------------------------

With the full LJ ball--surface interaction ($c_{ps}=1$,
$\eps_{ps}=\eps$), we run the same $N\in\{13,55,79,87\}$ at the 6 heights
($v_{\rm imp}=0.71$--$2.03\,\sigma/\tau$). The ball sticks in 24/24 cases
(Table~\ref{tab:adhesion}): no post-contact apex is detected, and
$Z_{\rm cm}$ at the end of the simulation ($400\,\tau$) sits at
$\approx10\,\sigma$, on top of the surface.

\begin{table*}[tp]
\centering
\caption{Adhesion energy and estimated threshold for $c_{ps}=1$,
$\eps_{ps}=\eps$.
$E_{\rm adh}=|\langle U_c\rangle_{\rm rest}|$ (average over the last 50
outputs of each simulation). $h_0^*$ estimated from the energy balance
$\alpha_{\rm cm}\,E_0^*=E_{\rm adh}$ with $\alpha_{\rm cm}\approx0.35$.}
\label{tab:adhesion}
\begin{tabular}{ccccc}
\toprule
$N$ & $E_{\rm adh}$ ($\eps$) & $E_{\rm adh}/N$ ($\eps$/atom) &
estimated $h_0^*$ ($\sigma$) & Scaling \\
\midrule
13 & 52  & 4.00 & 439 &  \\
55 & 135 & 2.45 & 270 &  \\
79 & 197 & 2.49 & 274 & $h_0^*\propto N^{-0.27}$ \\
87 & 209 & 2.40 & 264 &  \\
\bottomrule
\end{tabular}
\end{table*}

\paragraph{Adhesion energy.}
$E_{\rm adh}$ scales as $N^{0.73}$ (Figure~\ref{fig:adhesion}, right
panel), close to the contact-area scaling $\propto N^{2/3}=N^{0.67}$.
The energy per atom $E_{\rm adh}/N\approx2.4$--$4.0\,\eps$ decreases with
$N$, reflecting that interior atoms of the ball contribute less to the
contact than surface atoms.

\paragraph{Bounce threshold.}
The energy balance for the threshold requires the CM energy available after
the bounce to equal the adhesion energy:
\begin{equation}
    h_0^*(N) = \frac{E_{\rm adh}(N)}{\alpha_{\rm cm}\,Mg}
             = \frac{E_{\rm adh}(N)}{0.35\,N\gstar}.
    \label{eq:h0star}
\end{equation}
The estimated values ($h_0^*\approx260$--$440\,\sigma$,
Table~\ref{tab:adhesion}) lie outside the studied height range
($h_{0,\rm max}=79.5\,\sigma$). The dependence $h_0^*\propto N^{-0.27}$
indicates that \emph{larger} clusters have \emph{lower} thresholds: their
mass grows as $N$ but the adhesion only as $N^{0.73}$, making escape easier
for large projectiles.

This threshold should be read as an \emph{order-of-magnitude} bound, not as
a quantitative prediction, for three reasons. First, it is an extrapolation
to $3$--$5$ times the maximum validated height. Second, at those heights
($v_{\rm imp}\sim3.7$--$4.8\,\sigma/\tau$) the cluster no longer bounces as
a solid: \cite{Gutierrez2012} observed melting of the ball for
$h_0>41.5\,\sigma$, so in the threshold regime the contact area ---and with
it $E_{\rm adh}$--- would grow with $v$, instead of being constant as
Eq.~\eqref{eq:h0star} assumes. Third, $\alpha_{\rm cm}\approx0.35$ is taken
from the repulsive case; with adhesion the CM energy retention will be
smaller. All three biases point in the same direction (delaying the
bounce), so $h_0^*$ is likely a lower bound. Verifying this requires
$c_{ps}=1$ simulations at intermediate-to-high $h_0$ measuring
$E_{\rm adh}(v)$.

\begin{figure*}[tp]
    \centering
    \includegraphics[width=\textwidth]{../data/outputs/fig_paso_C_adhesion_en.pdf}
    \caption{\textbf{(a)}~Trajectories $Z_{\rm cm}(t)$ for $N=79$,
    $c_{ps}=1$, 6 heights (first $120\,\tau$; there are no further events
    afterwards). All end at $Z_{\rm cm}\approx10\,\sigma$ (surface top,
    dashed gray line): adhesion in all 6 cases.
    \textbf{(b)}~Threshold $h_0^*(N)$ estimated from
    Eq.~\eqref{eq:h0star} (red circles; the error bar propagates the
    spread of $E_{\rm adh}$ across heights), compared with the fit
    $E_{\rm adh}=7.9\,N^{0.73}$ (dashed black) and with the scaling
    $\propto N^{-1/3}$ (dotted blue).}
    \label{fig:adhesion}
\end{figure*}

% =============================================================================
\section{Discussion}
\label{sec:discusion}
% =============================================================================

\subsection{Origin of \texorpdfstring{$\alpha_{\Delta V_B}$}{α\_ΔVB}:
post-impact anharmonic expansion}
\label{sec:anarmonico}

The channel $\alpha_{\Delta V_B}$ has no analogue in rigid-sphere models.
During compression, the impulse from the surface converts CM KE into
intra-ball PE. Upon rebound, the CM recovers KE, but the relaxation of the
internal PE toward equilibrium takes several vibrations
($\sim1$--$5\,\tau$), by which time the cluster is already receding. At the
first apex, $V_B>V_B^{(0)}$, and the difference records the potential
energy the ball ``carries away''.

This asymmetry has two consequences: (i)~any partition using only the
internal KE $T_B^{(1)}$ underestimates by a factor $\sim2$ the heat
deposited in the ball; (ii)~the natural quantity for the equipartition
balance is $\alpha_{\rm ball}=\alpha_{T_B}+\alpha_{\Delta V_B}$, not
$\alpha_{T_B}$.

\subsection{Effective Nc and scaling with N}
\label{sec:Nc_interpretacion}

The calibration $N_c\approx N$ from $\alpha_{\rm ball}$ for $N\geq55$ has a
direct reading: the impact energy is distributed among the $3N-3$ internal
modes of the ball and $3N$ active surface modes, i.e., the active footprint
of the impact effectively involves $N$ surface atoms. This is consistent
with the ball radius ($\sim N^{1/3}$) also setting the penetration depth of
the pressure wave into the surface, activating $\sim N$ atoms in volume
(not only over the contact area $\sim N^{2/3}$).

For $N=13$, the system does not reach the equipartition regime: the ball is
too stiff (a single spherical shell) and $\alpha_{\rm ball}\approx0.11$
falls far below the predicted $\sim0.48$. The result marks a transition
between the ``molecular'' regime ($N\leq13$) and the ``bulk-like'' regime
($N\geq55$).

\subsection{Long-time dynamics and collapse}
\label{sec:largo_discusion}

The $5000\,\tau$ simulation reveals that $h^*=T_B/(Mg)$ is not a stationary
attractor but an approximate bound valid only for the first bounce (when
$\alpha_{\rm cm}\approx\alpha_{T_B}$). As $h_n$ decreases geometrically and
$T_B$ decreases more slowly, the ball approaches the surface until it
remains trapped above it. The collapse through the surface lattice
(detectable by $Z_{\rm cm}<z_0$) is an artifact of the purely repulsive
potential, which has no lateral barrier: with $c_{ps}=1$ this collapse does
not occur because the ball sticks before that stage.

\subsection{Adhesion threshold and soft landing}
\label{sec:adhesion_discusion}

The $c_{ps}=1$ results show that the LJ adhesion energy is considerably
larger than the available CM energy for all studied $N$ and $h_0$. The
threshold $h_0^*(N)\propto N^{-0.27}$ has a direct implication for
\emph{soft landing} applications: larger clusters require lower relative
impact velocities to stick ($v^*\propto N^{-0.13}$), which favors the
deposition of large clusters at intermediate velocities.

% =============================================================================
\section{Conclusions}
\label{sec:conclusiones}
% =============================================================================

We have reproduced and extended the study of Gutiérrez et
al.~\cite{Gutierrez2012} with \lammps{} using pre-minimized FCC clusters of
exact size ($N\in\{1,2,7,13,55,79,87\}$). The main conclusions are:

\begin{enumerate}

    \item \textbf{The impact involves four distinct energy channels.}
    $\alpha_{\rm cm}+\alpha_{T_B}+\alpha_{\Delta V_B}+\alpha_{\rm surf}=1$,
    with $\alpha_{T_B}\approx0.22$ constant,
    $\alpha_{\Delta V_B}\approx0.21$--$0.33$ and
    $\alpha_{\rm surf}\approx0.07$--$0.16$ for $N=79$,
    $v>0.9\,\sigma/\tau$. The anharmonic expansion dominates the direct
    surface excitation by a factor of $\sim3$--$4$.

    \item \textbf{The total internal fraction $\alpha_{\rm ball}$ scales
    according to equipartition with $N_c\approx N$ for $N\geq55$.}
    This predicts $\alpha_{\rm ball}\to1/2$ as $N\to\infty$, although the
    data systematically exceed that bound by $\sim8\%$
    ($N_c^{\rm free}\approx0.8$--$1.1\,N$; see Section~\ref{sec:sweep_N}).
    Using $\alpha_{T_B}$ alone, $N_c\propto N^{0.20}$ does not converge to
    a simple interpretation.

    \item \textbf{$N=13$ is in a different regime.}
    It does not reach the equipartition plateau within
    $v\leq2.0\,\sigma/\tau$: the stiffness of the smallest complete FCC
    shell prevents fast internal thermalization.

    \item \textbf{The scale $h_1\approx T_B/(Mg)$ is valid for the first
    bounce (ratio$=0.96$) but not as an attractor.}
    The CM height decays geometrically while $T_B$ decreases more slowly,
    and the ball remains trapped over the surface
    ($h_n\lesssim1\,\sigma$) after $\sim23$ bounces ($t\approx660\,\tau$
    for $h_0=28\,\sigma$, $N=79$). The subsequent leaking of atoms through
    the lattice ($t\approx713\,\tau$) is an artifact of the purely
    repulsive potential, not part of the trapping physics.

    \item \textbf{With adhesion ($c_{ps}=1$), the ball sticks in all
    studied cases ($v$ up to $2.0\,\sigma/\tau$).}
    $E_{\rm adh}\propto N^{0.73}$, $h_0^*(N)\propto N^{-0.27}$.
    The extrapolated threshold $h_0^*\sim260$--$440\,\sigma$ is an
    order-of-magnitude bound (likely a lower one; see
    Section~\ref{sec:adhesion}); the scaling favorable to large clusters
    has implications for \emph{soft landing}.

\end{enumerate}

As future work, the extension to initial temperature $\Tstar>0$ and to
multiple cluster orientations (3--5 replicas per point) will allow error
bars to be assigned to the fractions $\alpha_i$ ---and to decide whether
the non-monotonicity of $e_1$ with $h_0$ (Table~\ref{tab:particion4}) is
physical or initial-condition variance---, to quantify the statistical
dependence of $n_{\rm adh}$ (number of bounces before sticking), and to
assess the role of internal melting in the cascade. Measuring the
post-bounce angular momentum (\texttt{compute angmom}) will separate the
rotational contribution of $\alpha_{\rm ball}$, and sweeping the
intermediate complete shells $N=19$ and $43$ will populate the
molecular$\to$bulk transition between $N=13$ and $N=55$.

\begin{acknowledgments}
The authors thank the Master's program in Computational Simulation of the
PUCV for academic support. The simulations were performed with
\lammps~7~Feb~2024~\cite{Thompson2022}. The source code is available at
\url{https://github.com/fwasaff/inelastic-rebound-md}.
\end{acknowledgments}

\begin{thebibliography}{99}

\bibitem{Gutierrez2012}
G. Gutiérrez, S. Davis, C. Loyola, J. Peralta, F. González, Y. Navarrete
and F. González-Wasaff,
\textit{Inelastic Collisions and Hypervelocity Impacts at Nanoscopic Level:
A Molecular Dynamics Study},
in: \textit{Molecular Dynamics -- Theoretical Developments and Applications
in Nanotechnology and Energy}, chap.~12, pp.~229--248, InTech (2012).

\bibitem{Haff1983}
P.~K. Haff,
\textit{Grain flow as a fluid-mechanical phenomenon},
J. Fluid Mech. \textbf{134}, 401 (1983).

\bibitem{Dahneke1975}
B. Dahneke,
\textit{Further measurements of the bouncing of small latex spheres},
J. Colloid Interface Sci. \textbf{51}, 58 (1975).

\bibitem{Awasthi2007}
A. Awasthi, S.~C. Hendy, P. Zoontjens and S.~A. Brown,
\textit{Molecular dynamics simulations of reflection and adhesion behavior
in Lennard-Jones cluster deposition},
Phys. Rev. B \textbf{76}, 115437 (2007).

\bibitem{Awasthi2010}
A. Awasthi, S.~A. Brown and S.~C. Hendy,
\textit{Oblique impacts of Lennard-Jones clusters},
Phys. Rev. B \textbf{82}, 073406 (2010).

\bibitem{Northby1987}
J.~A. Northby,
\textit{Structure and binding of Lennard-Jones clusters: $13 \le N \le 147$},
J. Chem. Phys. \textbf{87}, 6166 (1987).

\bibitem{Wales1997}
D.~J. Wales and J.~P.~K. Doye,
\textit{Global optimization by basin-hopping and the lowest energy
structures of Lennard-Jones clusters containing up to 110 atoms},
J. Phys. Chem. A \textbf{101}, 5111 (1997).

\bibitem{Haberland1994}
H. Haberland, M. Mall, M. Moseler, Y. Qiang, T. Reiners and Y. Thurner,
\textit{Filling of micron-sized contact holes with copper by energetic
cluster impact},
J. Vac. Sci. Technol. A \textbf{12}, 2925 (1994).

\bibitem{Thompson2022}
A.~P. Thompson \textit{et al.},
\textit{LAMMPS -- A flexible simulation tool for particle-based materials
modeling},
Comput. Phys. Commun. \textbf{271}, 108171 (2022).

\end{thebibliography}

\end{document}
